% !TEX root = ../auv-thesis.tex

\begin{abstract}
  随着海上风电向深远海和规模化方向发展，海底电力电缆的敷设里程、作业水深与运维复杂度持续增加。传统依赖母船、潜水员或缆控水下机器人的巡检方式成本高、覆盖效率有限，难以满足长距离海缆的周期性前置筛查需求。自主式水下航行器（autonomous underwater vehicle, AUV）具备长航时、近底自主航行和多源传感能力，但水下声学、磁场、惯性与测速观测会随底质、掩埋状态、外部干扰和通信条件持续退化。若系统只传递状态估计值而不传递其可信程度，决策与控制层就可能把不可靠观测当作真值，进而放大跟踪误差和安全风险。因此，如何使航行器在感知质量变化时识别自身能力边界，并据此调整任务行为和控制强度，是海缆自主巡检由算法验证面向工程部署的关键问题。

  针对上述问题，本文提出一套以不确定性显式流转为主线的 AUV 海缆安全自主巡检系统。硬件上采用 Jetson Orin 感知决策脑与 PC104/VxWorks 安全执行脑构成异步双脑架构，将高算力非实时算法与硬实时执行保护解耦；软件上通过 ROS2、轻量 UDP 协议、上位机授权和底层看门狗形成分层安全闭环。系统由接口、感知、状态估计、决策和控制等功能层组成，并以统一时间戳、坐标系和置信度接口连接各层，使仿真、数字孪生与实物部署能够共享数据契约。为避免用单一指标替代工程证据，本文进一步建立由算法回归、故障注入、高保真数字孪生、部署接口闭环和实物待验层组成的验证阶梯，使每项结论均对应明确的证据等级。

  在感知与状态估计侧，本文以误差状态扩展卡尔曼滤波为骨架，融合惯性、多普勒测速、声学与磁场观测，并针对多速率、异步到达和观测质量变化设计时间对齐与自适应量测更新机制。磁学通道采用数字锁相放大提取工频幅值、相位和信噪比，利用三轴磁场模量的旋转不变性与之字形横切形成的半高全宽几何关系反演电缆垂直距离，从而降低对未知电流幅值、铠装衰减和传感器姿态的依赖。滤波器利用归一化新息平方和观测质量分数在线调节量测噪声协方差，再将状态协方差映射为归一化置信度，形成由观测质量、协方差与置信度构成的不确定性通道，为下游决策和控制提供统一输入。

  在决策与控制侧，本文采用行为树组织高优先级安全抢占，以任务状态机管理搜索、确认、跟踪和回收等生命周期阶段，并建立 ROS2 仲裁、自治守卫和 PC104/VxWorks 本地保护组成的分层安全逻辑。在控制层，将在线置信度引入模型预测控制的跟踪权重与控制代价，使观测可靠时优先提高跟踪精度、观测退化时主动降低控制激进程度，形成不确定性感知模型预测控制。该设计不假定单一控制器在所有路径和扰动下占优，而是把求解失败回退、安全约束和状态估计退化共同纳入闭环。

  实验结果表明：误差状态估计在八类传感器与通信扰动场景、每场景三个随机种子的二十四次运行中全部成功，自适应协方差机制在全部场景被触发，触发率为 0.24--0.36；磁外参仿真标定（单次标定）使平移误差和安装角误差分别降低 96.27\% 和 95.31\%。比例积分微分地形跟随在低、中、高三档地形中各完成三次运行，未出现低于安全阈值的违规，离底高度均方根误差为 0.59--0.71~m；在 low/mid/high 三档 CBF 60~s 三种子矩阵中，PID/MPC 共 18 次运行全部有效，最小净空为 2.70--2.90~m，1.5~m 低净空违规率和海床穿透率均为 0；MPC 求解 p95 上界为 23.84~ms，作为实时性边界保留。六类代理电缆场景的 R13-v2 正式矩阵完成 36/36 次运行，固定权重与保守分源置信控制的 fallback 均为 0，p95 求解耗时约 9.8--9.9~ms；保守分源置信控制未系统改善横向 RMSE，但将全局控制变化率 RMS 由 0.247 降至 0.085。综合极端 baseline full-flow 三种子补充实验全部完成，cross-track q95 均值为 0.206~m，fallback 与安全违规均为 0。端到端电缆探测链在确定性干净先验下三次数字孪生运行均就绪并通过；在经场景几何重配以恢复磁观测可观测性（将真值直缆重置于传感器正下方约 7.5 m）的前提下，中、重两档先验畸变的六自由度闭环恢复亦各完成三次窗口内验收。

  本文的主要贡献不在于宣称某一算法在所有工况下全面占优，而在于建立“状态与置信度共同输出、不确定性跨层传递、控制策略随可信度调整”的系统化技术路线，并以可追溯的分层证据给出其成立条件与失效边界。上述结果主要来自仿真、数字孪生与等效部署环境，可支撑算法正确性、软件接口和闭环数据契约，尚不能替代真实检测噪声、真机网络时延、硬件标定与海试条件下的实物验收；这一边界由本文有意采用的分层验证阶梯所界定，其向实物层推进的验证路径已作为后续工作明确给出。

  实物侧另以真实 PC104/VxWorks 验证零执行器通信与同步安全状态链：Bit5/Bit13/Bit14 已在同一 rosbag 时间轴驱动 ROS2 仲裁与行为树降级，修正后同步故障运行 18/18 项通过；独立 1800~s 稳态区间的 124841 个仲裁状态样本保持 100\% ACTIVE，目标故障状态样本为 0。因此，上述``真机网络时延尚不能替代''特指无共享时钟条件下不能报告严格单向物理时延，且现有实物结果仍不包含真实检测噪声、非零执行机构或水域验收。

  \thusetup{
    keywords = {自主式水下航行器, 海底电缆巡检, 声磁协同感知, 不确定性量化, 模型预测控制},
  }
\end{abstract}

\begin{abstract*}
  As offshore wind farms expand toward deeper and more remote waters, subsea power cables are being deployed over longer distances and in increasingly complex operating environments. Conventional inspection methods that rely on support vessels, divers, or tethered remotely operated vehicles are costly and inefficient for periodic, large-scale screening. Autonomous underwater vehicles (AUVs) provide long endurance, near-seabed autonomy, and multisensor payload capacity. However, acoustic, magnetic, inertial, and velocity observations may degrade continuously with seabed conditions, cable burial, electromagnetic interference, and communication quality. If an autonomous system propagates only state estimates without their confidence, its decision and control layers may treat unreliable observations as ground truth and amplify tracking errors and safety risks.

  This thesis develops a safe autonomous subsea-cable inspection system in which uncertainty is propagated explicitly across system layers. An asynchronous dual-computer architecture combines a Jetson Orin perception-and-decision computer with a PC104/VxWorks safety-execution computer. Robot Operating System 2 (ROS2), a lightweight User Datagram Protocol (UDP), operator authorization, and low-level watchdogs form a layered safety loop. For state estimation, an error-state extended Kalman filter fuses inertial, Doppler velocity, acoustic, and magnetic observations. A digital lock-in amplifier extracts power-frequency magnitude, phase, and signal-to-noise ratio. The rotation-invariant magnetic magnitude and the full-width-at-half-maximum geometry generated by a zigzag crossing are then used to infer the vertical cable distance with reduced dependence on unknown current magnitude, armoring attenuation, and sensor attitude. Normalized innovation squared statistics and observation-quality scores adapt the measurement-noise covariance online, while the estimated covariance is mapped to a confidence signal shared with downstream modules. For decision and control, a behavior tree provides safety preemption, a mission state machine manages the task life cycle, and the online confidence modulates the tracking weights and control penalties of model predictive control, yielding uncertainty-aware model predictive control. A layered validation framework is established to separate algorithm regression, fault injection, high-fidelity digital-twin evaluation, deployment-interface closed loops, and pending physical-vehicle verification.

  Across eight sensor and communication disturbance scenarios, the estimator completed all 24 runs with three random seeds per scenario. Covariance adaptation was activated in every scenario, with activation ratios between 0.24 and 0.36. Simulated magnetic-extrinsic calibration (a single calibration run) reduced the translation and installation-angle errors by 96.27\% and 95.31\%, respectively. The proportional--integral--derivative terrain-following controller completed three runs at each of three terrain intensities without violating the minimum-clearance threshold, and achieved clearance root-mean-square errors of 0.59--0.71~m. In the low/mid/high terrain CBF 60-s three-seed matrix, all 18 PID/MPC runs maintained a minimum clearance of 2.70--2.90~m, with zero 1.5-m clearance violations and zero seabed penetration; the MPC p95 solution time had an upper bound of 23.84~ms and is treated as a realtime boundary. The R13-v2 formal matrix over six proxy cable scenarios completed all 36 runs; both fixed weights and conservative source-specific uncertainty-aware control had zero fallback events, and the p95 solution time was approximately 9.8--9.9~ms. The conservative policy did not systematically improve lateral RMSE, but it reduced the global control-rate RMS from 0.247 to 0.085. A supplementary three-seed baseline full-flow run in the combined extreme scenario achieved a mean cross-track q95 of 0.206~m, with zero fallback and zero safety violations. The end-to-end cable-inspection chain was ready and passed in three digital-twin runs with a deterministic clean prior. After the scenario geometry was reconfigured to restore magnetic observability (relocating the ground-truth straight cable to about 7.5 m directly beneath the sensor), three runs at each of the medium and severe prior-distortion levels also passed the windowed acceptance criteria in a six-degree-of-freedom closed loop.

  The principal contribution is therefore not a claim that one algorithm dominates under all conditions. Instead, the thesis establishes a system-level route in which state and confidence are produced together, uncertainty is transferred across layers, and control behavior is adjusted according to confidence, with traceable evidence defining where the route is effective. The reported results are primarily obtained in simulation, digital-twin, and deployment-equivalent environments. They support algorithmic correctness, software interfaces, and closed-loop data contracts, but do not replace physical acceptance with realistic detection noise, hardware calibration, real network latency, and sea trials. This boundary is defined by the layered validation ladder adopted deliberately in this thesis, whose path toward physical-layer verification is stated explicitly as future work.

  On the physical side, a real PC104/VxWorks computer was used to verify the zero-actuator communication and synchronized safety-state chain. Fault Bits 5, 13, and 14 drove ROS2 arbitration and behavior-tree degradation on one rosbag timeline, and the corrected run passed all 18 checks. In a separate 1800-s steady-state interval, all 124841 arbiter-status samples remained ACTIVE and no target fault state occurred. The remaining network-timing limitation refers specifically to strict one-way physical latency without a shared clock; realistic detection noise, nonzero actuator dynamics, and water trials also remain outside the present evidence.

  \thusetup{
    keywords* = {autonomous underwater vehicle, subsea cable inspection, acoustic-magnetic cooperative perception, uncertainty quantification, model predictive control},
  }
\end{abstract*}
